\documentclass[11pt]{article}

\usepackage[a4paper,margin=1in]{geometry}
\usepackage{amsmath,amssymb,amsthm,mathtools}
\usepackage{enumitem}
\usepackage{hyperref}
\usepackage[nameinlink,noabbrev]{cleveref}

\hypersetup{
  colorlinks=true,
  linkcolor=blue,
  citecolor=blue,
  urlcolor=blue
}

\newtheorem{theorem}{Theorem}[section]
\newtheorem{lemma}[theorem]{Lemma}
\newtheorem{proposition}[theorem]{Proposition}
\newtheorem{corollary}[theorem]{Corollary}
\theoremstyle{definition}
\newtheorem{definition}[theorem]{Definition}
\newtheorem{remark}[theorem]{Remark}
\newtheorem{convention}[theorem]{Convention}

\newcommand{\Ind}{\operatorname{Ind}}
\newcommand{\Res}{\operatorname{Res}}
\newcommand{\Aut}{\operatorname{Aut}}
\newcommand{\End}{\operatorname{End}}
\newcommand{\Hom}{\operatorname{Hom}}
\newcommand{\Stab}{\operatorname{Stab}}
\newcommand{\op}{\mathrm{op}}
\newcommand{\Mod}{\mathrm{Mod}}
\newcommand{\Span}{\operatorname{span}}
\newcommand{\supp}{\operatorname{supp}}

\title{A Formalization-Oriented Blueprint for Mackey's Formula for Representations}
\author{}
\date{}

\begin{document}

\maketitle

\begin{abstract}
This note gives a detailed mathematical proof of Mackey's restriction formula for induced representations in a form intended to guide a future formalization using the representation-theoretic \texttt{Representation} API. No Lean code is included. The proof is written in a generator-and-relation style compatible with a common construction of induced representations as a quotient or coinvariant module. The final section specializes the formula to the index-two subgroup statement: if $K \leq G$ has index $2$, $\psi : K \to k^\times$ is a character, $\rho = \Ind_K^G(\psi)$, and $J \leq G$ is not contained in $K$, then
\[
  \Res^G_J \rho \cong \Ind_{K \cap J}^{J}(\psi|_{K\cap J}).
\]
The equality should be formalized as an isomorphism of $J$-representations, not as definitional equality.
\end{abstract}

\tableofcontents

\section{Purpose and ambient assumptions}

The goal is to formalize the Mackey formula for finite groups and finite-index subgroups, using the usual representation-theoretic induction functor.

\begin{convention}[Algebraic setting]
Let $k$ be a commutative ring. In applications to characters, one may take $k$ to be an algebraically closed field, but the Mackey formula itself does not require algebraic closedness.

Let $G$ be a finite group, and let $H,L \leq G$ be subgroups. Let $V$ be a $k$-module and let
\[
  \rho : H \to \Aut_k(V)
\]
be a $k$-linear representation of $H$ on $V$.
\end{convention}

\begin{remark}[Expected Mathlib background]
For a Mathlib development, one may assume or import the following standard ingredients.
\begin{enumerate}[label=(\alph*)]
  \item The type of representations, usually represented by a monoid or group homomorphism into $k$-linear endomorphisms of a module.
  \item Restriction of a representation along a group homomorphism or subgroup inclusion.
  \item Induction of representations along a group homomorphism or subgroup inclusion.
  \item Basic finite group and subgroup facts: cosets, quotient by a subgroup, double cosets, finite direct sums, and finite support functions.
  \item The theorem that a subgroup of index $2$ is normal. This is useful background, though the index-two corollary below can be proved directly from the coset count.
\end{enumerate}
The proof below deliberately avoids Lean identifiers and gives only the mathematical structure that should be formalized.
\end{remark}

\section{Conventions for induced representations}

The induced representation $\Ind_H^G(V)$ will be described by generators
\[
  [g,v], \qquad g\in G,\ v\in V,
\]
subject to the relation
\begin{equation}\label{eq:ind-relation}
  [hg,\rho(h)v] = [g,v]
  \qquad (h\in H,\ g\in G,\ v\in V).
\end{equation}
Equivalently,
\begin{equation}\label{eq:ind-relation-inverse}
  [hg,v] = [g,\rho(h)^{-1}v].
\end{equation}
The induced $G$-action is taken to be the right-translation convention
\begin{equation}\label{eq:G-action}
  x\cdot[g,v] = [gx^{-1},v]
  \qquad (x,g\in G,\ v\in V).
\end{equation}
Consequently, the restriction to $L\leq G$ acts by
\begin{equation}\label{eq:L-action-restricted}
  \ell\cdot[g,v]=[g\ell^{-1},v]
  \qquad (\ell\in L).
\end{equation}

\begin{remark}[Why this convention]
This convention matches the standard proof of Mackey's formula in which $\Ind_H^G(V)$ decomposes as a sum over left cosets $H\backslash G$, and the restricted $L$-action permutes those summands by right multiplication. With this convention, the subgroup attached to a double coset representative $s$ is
\[
  H_s := s^{-1}Hs\cap L,
\]
and the twisted representation is
\[
  \rho^s(a)=\rho(sas^{-1}).
\]
\end{remark}

\section{The left-coset decomposition of induction}

Let $T\subseteq G$ be a choice of representatives for the left cosets $H\backslash G$. Thus each $g\in G$ has a unique expression
\[
  g=ht
\]
with $h\in H$ and $t\in T$.

For each $t\in T$, let $V_t$ denote a copy of $V$. Define a $k$-linear map
\[
  \Phi_T : \bigoplus_{t\in T} V_t \longrightarrow \Ind_H^G(V)
\]
by sending $v\in V_t$ to
\[
  [t,v].
\]

\begin{proposition}[Induction is a sum over left cosets]\label{prop:left-coset-decomp}
The map
\[
  \Phi_T : \bigoplus_{t\in T} V_t \longrightarrow \Ind_H^G(V)
\]
is a $k$-linear isomorphism.
\end{proposition}

\begin{proof}
First prove surjectivity. Let $[g,v]$ be a generator of $\Ind_H^G(V)$. Write $g=ht$ with $h\in H$ and $t\in T$. Then by \eqref{eq:ind-relation-inverse},
\[
  [g,v]=[ht,v]=[t,\rho(h)^{-1}v].
\]
Thus every generator is in the image of the $t$-summand map, so $\Phi_T$ is surjective.

To prove injectivity, define an inverse map on generators. Given $g\in G$, write $g=ht$ uniquely with $h\in H$ and $t\in T$, and set
\[
  \Psi_T([g,v])
  :=
  \text{the element of }\bigoplus_{t'\in T}V_{t'}
  \text{ supported at }t\text{ with value }\rho(h)^{-1}v.
\]
We must verify that this respects the induction relation. Let $a\in H$. If $g=ht$, then
\[
  ag=(ah)t.
\]
Therefore
\[
  \Psi_T([ag,\rho(a)v])
  =
  \text{the vector supported at }t\text{ with value }\rho(ah)^{-1}\rho(a)v.
\]
Since $\rho(ah)=\rho(a)\rho(h)$, we have
\[
  \rho(ah)^{-1}\rho(a)v
  =
  \rho(h)^{-1}\rho(a)^{-1}\rho(a)v
  =
  \rho(h)^{-1}v.
\]
Hence
\[
  \Psi_T([ag,\rho(a)v])=\Psi_T([g,v]).
\]
So $\Psi_T$ is well-defined.

The composites are immediate on generators. For $v\in V_t$,
\[
  \Psi_T(\Phi_T(v))=\Psi_T([t,v])=v
\]
because $t=1\cdot t$. Conversely, if $g=ht$, then
\[
  \Phi_T(\Psi_T([g,v]))
  =
  [t,\rho(h)^{-1}v]
  =
  [ht,v]
  =
  [g,v].
\]
Thus $\Phi_T$ is an isomorphism.
\end{proof}

\begin{remark}[Formalization target]
The proof of \cref{prop:left-coset-decomp} is the basic normal-form lemma for induced representations. In a formal development, it is useful to state this both for a chosen transversal and invariantly as a direct sum indexed by the quotient type $H\backslash G$.
\end{remark}

\section{Restriction to a subgroup and double-coset decomposition}

The subgroup $L$ acts on the set of left cosets $H\backslash G$ by right multiplication:
\[
  (Hg)\cdot \ell := H(g\ell^{-1}).
\]
This action is well-defined because if $Hg=Hg'$, then $g'=hg$ for some $h\in H$, so
\[
  g'\ell^{-1}=hg\ell^{-1},
\]
and therefore
\[
  H(g'\ell^{-1})=H(g\ell^{-1}).
\]
The inverse in the formula appears because the representation action is $\ell\cdot[g,v]=[g\ell^{-1},v]$.

The orbits of this right $L$-action on $H\backslash G$ are precisely the double cosets $H\backslash G/L$. More explicitly, the orbit of $Hs$ is
\[
  \{Hs\ell : \ell\in L\},
\]
which corresponds to the double coset $HsL$.

For a double coset $D=HsL$, define
\[
  M_D := \Span_k\{[g,v] : g\in D,\ v\in V\}\subseteq \Ind_H^G(V).
\]
Equivalently, after choosing a left-coset decomposition, $M_D$ is the direct sum of all summands indexed by left cosets contained in $D$:
\[
  M_D = \bigoplus_{Hg\subseteq D} V_{Hg}.
\]

\begin{lemma}[Double-coset summands are stable]\label{lem:double-coset-stable}
For each double coset $D=HsL$, the submodule $M_D$ is stable under the restricted $L$-action.
\end{lemma}

\begin{proof}
It suffices to check generators. Let $g\in D$ and $\ell\in L$. Then $g\ell^{-1}\in DL=D$, because $D=HsL$ is stable under right multiplication by elements of $L$. Hence
\[
  \ell\cdot[g,v]=[g\ell^{-1},v]\in M_D.
\]
Thus $M_D$ is $L$-stable.
\end{proof}

\begin{proposition}[Restriction decomposes over double cosets]\label{prop:double-coset-decomp}
There is an isomorphism of $L$-representations
\[
  \Res^G_L\Ind_H^G(V)
  \cong
  \bigoplus_{D\in H\backslash G/L} M_D.
\]
\end{proposition}

\begin{proof}
As a $k$-module, $\Ind_H^G(V)$ decomposes as a direct sum over left cosets $H\backslash G$ by \cref{prop:left-coset-decomp}. The double cosets partition the set of left cosets. Therefore the direct sum over all left cosets can be regrouped as a direct sum over double cosets:
\[
  \bigoplus_{Hg\in H\backslash G}V_{Hg}
  \cong
  \bigoplus_{D\in H\backslash G/L}\left(\bigoplus_{Hg\subseteq D}V_{Hg}\right).
\]
By \cref{lem:double-coset-stable}, each inner summand is stable under the $L$-action. The regrouping map is therefore $L$-equivariant, and hence gives an isomorphism of $L$-representations.
\end{proof}

\section{The subgroup and twisted representation attached to a double coset}

Fix $s\in G$. Define
\[
  H_s := s^{-1}Hs\cap L.
\]
This is a subgroup of $L$.

Define a representation
\[
  \rho^s : H_s \to \Aut_k(V)
\]
by
\[
  \rho^s(a)=\rho(sas^{-1}).
\]
This is well-defined because $a\in H_s$ implies $a\in s^{-1}Hs$, hence $sas^{-1}\in H$.

\begin{lemma}[The twisted action is a representation]\label{lem:twisted-rep}
The map $\rho^s$ is a $k$-linear representation of $H_s$ on $V$.
\end{lemma}

\begin{proof}
For $a,b\in H_s$,
\[
  \rho^s(ab)
  =
  \rho(sabs^{-1})
  =
  \rho((sas^{-1})(sbs^{-1}))
  =
  \rho(sas^{-1})\rho(sbs^{-1})
  =
  \rho^s(a)\rho^s(b).
\]
Also
\[
  \rho^s(1)=\rho(1)=1.
\]
Thus $\rho^s$ is a representation.
\end{proof}

\section{Identifying one double-coset summand}

Let $D=HsL$. We now identify $M_D$ with an induced representation from $H_s$ to $L$.

Consider
\[
  \Ind_{H_s}^{L}(V^s),
\]
where $V^s$ denotes $V$ equipped with the $H_s$-action $\rho^s$. Its generators will be written
\[
  [\ell,v]_s,
  \qquad \ell\in L,
  \ v\in V,
\]
subject to
\begin{equation}\label{eq:Hs-ind-relation}
  [a\ell,\rho^s(a)v]_s=[\ell,v]_s
  \qquad (a\in H_s).
\end{equation}
The $L$-action is
\begin{equation}\label{eq:ind-L-action}
  \lambda\cdot[\ell,v]_s=[\ell\lambda^{-1},v]_s.
\end{equation}

Define a map on generators
\[
  \Theta_s : \Ind_{H_s}^{L}(V^s) \longrightarrow M_D
\]
by
\begin{equation}\label{eq:theta-def}
  \Theta_s([\ell,v]_s)=[s\ell,v].
\end{equation}

\begin{lemma}[Well-definedness of $\Theta_s$]\label{lem:theta-well-defined}
The map $\Theta_s$ is well-defined and $k$-linear.
\end{lemma}

\begin{proof}
The only nontrivial point is compatibility with the relation \eqref{eq:Hs-ind-relation}. Let $a\in H_s$. Then $sas^{-1}\in H$. Therefore
\[
  \Theta_s([a\ell,\rho^s(a)v]_s)
  =
  [sa\ell,\rho(sas^{-1})v].
\]
But
\[
  sa\ell=(sas^{-1})s\ell.
\]
Using the induction relation in $\Ind_H^G(V)$, we obtain
\[
  [(sas^{-1})s\ell,\rho(sas^{-1})v]
  =
  [s\ell,v].
\]
Thus
\[
  \Theta_s([a\ell,\rho^s(a)v]_s)
  =
  \Theta_s([\ell,v]_s).
\]
Hence $\Theta_s$ descends to the quotient defining induction, and so is a well-defined $k$-linear map.
\end{proof}

\begin{lemma}[$L$-equivariance of $\Theta_s$]\label{lem:theta-equivariant}
The map $\Theta_s$ is $L$-equivariant.
\end{lemma}

\begin{proof}
Let $\lambda,\ell\in L$ and $v\in V$. Then by \eqref{eq:ind-L-action} and \eqref{eq:theta-def},
\[
  \Theta_s(\lambda\cdot[\ell,v]_s)
  =
  \Theta_s([\ell\lambda^{-1},v]_s)
  =
  [s\ell\lambda^{-1},v].
\]
On the other hand, by the restricted action \eqref{eq:L-action-restricted},
\[
  \lambda\cdot\Theta_s([\ell,v]_s)
  =
  \lambda\cdot[s\ell,v]
  =
  [s\ell\lambda^{-1},v].
\]
The two expressions agree on generators, so $\Theta_s$ is $L$-equivariant.
\end{proof}

\begin{lemma}[Surjectivity of $\Theta_s$]\label{lem:theta-surj}
The map $\Theta_s$ is surjective onto $M_D$.
\end{lemma}

\begin{proof}
The submodule $M_D$ is generated by all $[g,v]$ with $g\in HsL$. Write
\[
  g=hs\ell
\]
with $h\in H$ and $\ell\in L$. Then
\[
  [g,v]
  =
  [hs\ell,v]
  =
  [s\ell,\rho(h)^{-1}v]
  =
  \Theta_s([\ell,\rho(h)^{-1}v]_s).
\]
Thus every generator of $M_D$ lies in the image.
\end{proof}

\begin{lemma}[Injectivity of $\Theta_s$]\label{lem:theta-inj}
The map $\Theta_s$ is injective.
\end{lemma}

\begin{proof}
Construct an inverse map
\[
  \Psi_s : M_D\longrightarrow \Ind_{H_s}^{L}(V^s).
\]
For a generator $[g,v]$ with $g\in HsL$, choose $h\in H$ and $\ell\in L$ such that
\[
  g=hs\ell.
\]
Define
\begin{equation}\label{eq:psi-def}
  \Psi_s([g,v]) := [\ell,\rho(h)^{-1}v]_s.
\end{equation}
We must check that this is independent of the chosen expression $g=hs\ell$.

Suppose
\[
  hs\ell=h's\ell'
\]
with $h,h'\in H$ and $\ell,\ell'\in L$. Then
\[
  h'^{-1}hs=s\ell'\ell^{-1},
\]
and hence
\[
  s^{-1}h'^{-1}hs=\ell'\ell^{-1}.
\]
Set
\[
  a:=\ell'\ell^{-1}.
\]
Then $a\in L$, and also
\[
  a=s^{-1}h'^{-1}hs\in s^{-1}Hs.
\]
Thus
\[
  a\in H_s=s^{-1}Hs\cap L.
\]
Moreover,
\[
  \ell'=a\ell.
\]

Now compute
\[
  \rho^s(a)
  =
  \rho(sas^{-1})
  =
  \rho(h'^{-1}h).
\]
Therefore
\[
  \rho^s(a)\rho(h)^{-1}v
  =
  \rho(h'^{-1}h)\rho(h)^{-1}v
  =
  \rho(h')^{-1}\rho(h)\rho(h)^{-1}v
  =
  \rho(h')^{-1}v.
\]
Using the relation in $\Ind_{H_s}^{L}(V^s)$,
\[
  [\ell',\rho(h')^{-1}v]_s
  =
  [a\ell,\rho^s(a)\rho(h)^{-1}v]_s
  =
  [\ell,\rho(h)^{-1}v]_s.
\]
Thus \eqref{eq:psi-def} is independent of the choice of $h$ and $\ell$.

Next, check that $\Psi_s$ respects the original induction relation defining $\Ind_H^G(V)$. Let $b\in H$. If $g=hs\ell$, then
\[
  bg=(bh)s\ell.
\]
So
\[
  \Psi_s([bg,\rho(b)v])
  =
  [\ell,\rho(bh)^{-1}\rho(b)v]_s.
\]
Since $\rho(bh)=\rho(b)\rho(h)$, this equals
\[
  [\ell,\rho(h)^{-1}v]_s
  =
  \Psi_s([g,v]).
\]
Hence $\Psi_s$ is well-defined on $M_D$.

Finally, the two composites are identities on generators. First,
\[
  \Psi_s(\Theta_s([\ell,v]_s))
  =
  \Psi_s([s\ell,v])
  =
  [\ell,v]_s,
\]
using the expression $s\ell=1\cdot s\ell$. Second, if $g=hs\ell$, then
\[
  \Theta_s(\Psi_s([g,v]))
  =
  \Theta_s([\ell,\rho(h)^{-1}v]_s)
  =
  [s\ell,\rho(h)^{-1}v]
  =
  [hs\ell,v]
  =
  [g,v].
\]
Thus $\Psi_s$ is the inverse of $\Theta_s$, so $\Theta_s$ is injective.
\end{proof}

\begin{proposition}[One double-coset summand]\label{prop:one-summand}
For each $s\in G$, there is an isomorphism of $L$-representations
\[
  \Ind_{s^{-1}Hs\cap L}^{L}(V^s)
  \cong
  M_{HsL}.
\]
\end{proposition}

\begin{proof}
The isomorphism is $\Theta_s$. It is well-defined by \cref{lem:theta-well-defined}, $L$-equivariant by \cref{lem:theta-equivariant}, surjective by \cref{lem:theta-surj}, and injective by \cref{lem:theta-inj}.
\end{proof}

\section{Mackey's restriction formula}

Choose a set $S\subseteq G$ of representatives for the double cosets $H\backslash G/L$.

\begin{theorem}[Mackey formula]\label{thm:mackey}
There is an isomorphism of $L$-representations
\[
  \Res^G_L\Ind_H^G(V)
  \cong
  \bigoplus_{s\in S}
  \Ind_{s^{-1}Hs\cap L}^{L}(V^s),
\]
where $V^s$ denotes $V$ with the action
\[
  a\cdot_s v := \rho(sas^{-1})v
  \qquad (a\in s^{-1}Hs\cap L).
\]
\end{theorem}

\begin{proof}
By \cref{prop:double-coset-decomp},
\[
  \Res^G_L\Ind_H^G(V)
  \cong
  \bigoplus_{D\in H\backslash G/L}M_D.
\]
Choosing a representative $s\in S$ for each double coset identifies $D$ with $HsL$. By \cref{prop:one-summand},
\[
  M_{HsL}
  \cong
  \Ind_{s^{-1}Hs\cap L}^{L}(V^s).
\]
Taking the direct sum of these isomorphisms over all $s\in S$ yields
\[
  \Res^G_L\Ind_H^G(V)
  \cong
  \bigoplus_{s\in S}
  \Ind_{s^{-1}Hs\cap L}^{L}(V^s).
\]
The resulting map is an isomorphism of $L$-representations because each summand map is $L$-equivariant.
\end{proof}

\begin{remark}[Representative dependence]
The displayed formula uses a chosen set $S$ of double-coset representatives. A more invariant version is indexed by the quotient type $H\backslash G/L$, with each orbit represented by a chosen representative. For formalization, it may be simpler to first state the theorem with an explicit representative function.
\end{remark}

\section{The index-two corollary}

Now specialize to the statement motivating the formalization.

\begin{convention}[Character as a one-dimensional representation]
Let $k$ be a field, and let
\[
  \psi : K \to k^\times
\]
be a character of a subgroup $K\leq G$. We view $k$ as a one-dimensional $K$-representation by
\[
  x\cdot a := \psi(x)a.
\]
Let
\[
  \rho := \Ind_K^G(\psi).
\]
More explicitly, $\rho$ is $\Ind_K^G(k_\psi)$, where $k_\psi$ is the one-dimensional representation associated to $\psi$.
\end{convention}

\begin{lemma}[The double-coset space is a singleton]\label{lem:index-two-singleton}
Let $K\leq G$ have index $2$, and let $J\leq G$ be a subgroup not contained in $K$. Then
\[
  K\backslash G/J
\]
has exactly one element. Equivalently,
\[
  KJ=G.
\]
\end{lemma}

\begin{proof}
Since $J$ is not contained in $K$, choose
\[
  j\in J\setminus K.
\]
Because $K$ has index $2$, there are exactly two right cosets of $K$ in $G$. They are
\[
  K
  \qquad\text{and}\qquad
  Kj.
\]
Indeed, $j\notin K$, so $Kj\neq K$, and hence $K$ and $Kj$ exhaust the two right cosets.

Let $g\in G$. If $g\in K$, then $g\in KJ$ because $1\in J$. If $g\notin K$, then $Kg=Kj$, so $g=kj$ for some $k\in K$, and again $g\in KJ$. Therefore
\[
  G\subseteq KJ.
\]
The reverse inclusion $KJ\subseteq G$ is automatic, so $KJ=G$.

It follows that the only double coset is
\[
  KJ=G.
\]
Thus $K\backslash G/J$ is a singleton.
\end{proof}

\begin{theorem}[Index-two Mackey corollary]\label{thm:index-two-corollary}
Let $K\leq G$ have index $2$. Let $\psi:K\to k^\times$ be a character, and let
\[
  \rho:=\Ind_K^G(\psi).
\]
If $J\leq G$ is not contained in $K$, then there is an isomorphism of $J$-representations
\[
  \Res^G_J \rho
  \cong
  \Ind_{K\cap J}^{J}(\psi|_{K\cap J}).
\]
\end{theorem}

\begin{proof}
Apply \cref{thm:mackey} with $H=K$ and $L=J$. By \cref{lem:index-two-singleton}, the double-coset set $K\backslash G/J$ has a single representative, which we may take to be $s=1$.

The corresponding subgroup is
\[
  1^{-1}K1\cap J=K\cap J.
\]
The corresponding twisted representation is
\[
  \psi^1(x)=\psi(1x1^{-1})=\psi(x).
\]
Therefore the single Mackey summand is
\[
  \Ind_{K\cap J}^{J}(\psi|_{K\cap J}).
\]
Since there is only one summand, Mackey's formula gives
\[
  \Res^G_J\Ind_K^G(\psi)
  \cong
  \Ind_{K\cap J}^{J}(\psi|_{K\cap J}).
\]
This is the desired isomorphism.
\end{proof}

\begin{remark}[Correction to informal notation]
Since $\psi$ is defined on $K$, the restriction appearing on the right-hand side should be written
\[
  \psi|_{K\cap J},
\]
not literally $\psi|_J$, unless one has first extended $\psi$ to a function on $J$.
\end{remark}

\section{Formalization roadmap, still without Lean code}

A useful order of formalization is the following.

\begin{enumerate}[label=\textbf{Step \arabic*.},wide]
  \item \textbf{Fix conventions for induction.}
  Confirm the precise relation used by the chosen induced representation construction. The proof above assumes the convention
  \[
    [hg,\rho(h)v]=[g,v]
  \]
  and the action
  \[
    x\cdot[g,v]=[gx^{-1},v].
  \]
  If the library convention differs by inverses, adjust all formulas accordingly.

  \item \textbf{Prove the left-coset normal form.}
  For a chosen transversal $T$ of $H\backslash G$, prove that every generator $[g,v]$ has normal form $[t,v']$. Then build the isomorphism
  \[
    \Ind_H^G(V)\cong \bigoplus_{t\in T}V.
  \]

  \item \textbf{Make the coset-indexed decomposition intrinsic.}
  Replace the transversal $T$ with the quotient type $H\backslash G$. This avoids carrying a concrete representative set through later statements.

  \item \textbf{Define double-coset summands.}
  For each double coset $D=HsL$, define the submodule $M_D$ generated by $[g,v]$ with $g\in D$. Prove $M_D$ is stable under $L$.

  \item \textbf{Regroup the direct sum.}
  Prove
  \[
    \Res^G_L\Ind_H^G(V)
    \cong
    \bigoplus_{D\in H\backslash G/L}M_D.
  \]

  \item \textbf{Define the twisted representation.}
  For $s\in G$, define
  \[
    H_s=s^{-1}Hs\cap L
  \]
  and
  \[
    \rho^s(a)=\rho(sas^{-1}).
  \]
  Prove this is a representation of $H_s$.

  \item \textbf{Identify one summand.}
  Define the map
  \[
    \Theta_s([\ell,v]_s)=[s\ell,v].
  \]
  Prove that it is well-defined, $L$-equivariant, surjective, and injective. The inverse is given by
  \[
    [hs\ell,v]\mapsto [\ell,\rho(h)^{-1}v]_s.
  \]

  \item \textbf{Assemble Mackey.}
  Sum the isomorphisms over chosen double-coset representatives.

  \item \textbf{Specialize to index two.}
  Prove that if $[G:K]=2$ and $J\not\leq K$, then $KJ=G$, so the double-coset quotient $K\backslash G/J$ is a singleton.
\end{enumerate}

\section{References}

The proof follows the standard finite-group Mackey decomposition for induced representations. Useful references include:

\begin{enumerate}
  \item J.-P. Serre, \emph{Linear Representations of Finite Groups}, especially the discussion of induced representations and Mackey's formula.
  \item B. Conrad, \emph{Mackey theory}, course handout, Stanford University. This source uses the double-coset proof with the formula
  \[
    \Res_K^G\Ind_H^G(\rho)
    \simeq
    \bigoplus_{s\in H\backslash G/K}
    \Ind_{s^{-1}Hs\cap K}^{K}(\rho^s).
  \]
  \item C. W. Curtis and I. Reiner, \emph{Methods of Representation Theory}, Vol. I, for a broader treatment of induction, restriction, and Mackey decomposition.
  \item Mathlib documentation for representation theory, especially the files concerning \texttt{Representation}, restriction, and induced representations.
\end{enumerate}

\end{document}
